\section{Limitations and conclusion}
\label{sec:conclusion}

\paragraph{Limitations.} Apart from one roadside estimate at Alfonso
Reyes, there are no traffic counts in this study. Demand is a parametric model, varied over wide ranges, and every absolute
delay should be read in that light. The conclusions we draw are the ones
that survive the whole Monte Carlo: where the queue forms, and which
predictors find it. The repository accepts field counts
(\texttt{data/field\_counts.csv}, vehicles per hour per approach and
direction, or vehicles per minute on green and red); the Alfonso Reyes
estimate shows how much one count can change: it removed the predicted
westbound queue. Fifteen minutes of video at a few junctions, especially
eastbound, would replace the largest remaining assumptions with data. OpenStreetMap is incomplete in two places
that matter. Signals are under-mapped, so we infer them at major
frontage junctions. The five-to-three-lane merge that dominates the
results is a mapped feature that should be checked on site; if the
frontage road actually continues there, the bottleneck moves, as the
filled-gap experiment shows. The CTM aggregates express and frontage
lanes into one cross-section, applies green ratios as average capacities
rather than cycles, treats each approach as a vertical queue, and gives
every vehicle the same route behaviour. How strongly the queue blocks
exits ($\phi$) changes delay several-fold and is the first thing a field
study should measure. The model has no route choice, so traffic does not
shift to Gómez Mor\'in, Vasconcelos or Garza Sada when L\'azaro C\'ardenas
fills up.

\paragraph{Conclusion.} The lane count behind the popular explanation is
right: about \InvFeederLanesNoService\ street lanes, \InvFeederLanes\
with driveways, into \InvMainlineLanes\ per direction. But it is the
wrong quantity. Feeder lanes carry a few percent of their saturation flow,
and their number does not locate the queue: the local lane ratio found
the bottleneck in \McHitLR\% of congested runs. Utilisation found it in
\McHitRho\% and predicted whether congestion occurs at all. Remove the
lane drop, though, and the lane ratio becomes the better locator
(\McGapHitLR\% against \McGapHitRho\%, in a small sample): where many
accesses join in a short stretch, their friction matters. On this
corridor, at this map, the answer is geometric: a merge from five lanes
to three between Alfonso Reyes and the Arroyo Seco bridge. The
intervention that helps in every run is capacity at that merge, not
fewer driveways or retimed signals. The rule generalises beyond
Monterrey: to find where a corridor jams, compare demand with capacity at
each merge, and look first for the place where lanes disappear.
